\section{Introduction}\label{sec:intro}

Start with an $h\times m$ grid and black out some of its points.
The game master chooses a rectangle (tilted ones are allowed) and tells
you which of its corners are white. Can you recover all four corners?
Kagey~\cite{kagey001,kageyG4G14} asks how many points can be blacked out
while keeping the answer unique for every rectangle.
Such a set is a \emph{valid blackout} (Definition~\ref{def:valid}).
One round is shown in Figure~\ref{fig:game}.

\begin{figure}[t]
\centering
\resizebox{0.8\textwidth}{!}{\begin{tikzpicture}[x=0.9cm,y=0.9cm,baseline=(current bounding box.north)]
  \begin{scope}[xshift=0.000cm]
    \draw[black,line width=0.36pt,dash pattern=on 1.35pt off 1.35pt] (1,-1) -- (4,-1) (1,-2) -- (4,-2) (1,-3) -- (4,-3) (1,-1) -- (1,-3) (2,-1) -- (2,-3) (3,-1) -- (3,-3) (4,-1) -- (4,-3);
    \filldraw[fill=white,draw=black,line width=0.54pt] (1,-1) circle (0.3);
    \filldraw[fill=white,draw=black,line width=0.54pt] (2,-1) circle (0.3);
    \filldraw[fill=white,draw=black,line width=0.54pt] (3,-1) circle (0.3);
    \filldraw[fill=white,draw=black,line width=0.54pt] (4,-1) circle (0.3);
    \filldraw[fill=white,draw=black,line width=0.54pt] (1,-2) circle (0.3);
    \filldraw[fill=white,draw=black,line width=0.54pt] (2,-2) circle (0.3);
    \filldraw[fill=kyellow,draw=black,line width=0.54pt] (3,-2) circle (0.3);
    \filldraw[fill=white,draw=black,line width=0.54pt] (4,-2) circle (0.3);
    \filldraw[fill=white,draw=black,line width=0.54pt] (1,-3) circle (0.3);
    \filldraw[fill=white,draw=black,line width=0.54pt] (2,-3) circle (0.3);
    \filldraw[fill=white,draw=black,line width=0.54pt] (3,-3) circle (0.3);
    \filldraw[fill=white,draw=black,line width=0.54pt] (4,-3) circle (0.3);
    \node[lab] at (1,-0.45) {1};
    \node[lab] at (2,-0.45) {2};
    \node[lab] at (3,-0.45) {3};
    \node[lab] at (4,-0.45) {4};
    \node[lab] at (0.45,-1) {1};
    \node[lab] at (0.45,-2) {2};
    \node[lab] at (0.45,-3) {3};
    \node[capt,anchor=north] at (2.50,-3.45) {(a) the point $(3,2)$};
  \end{scope}
  \begin{scope}[xshift=4.680cm]
    \draw[black,line width=0.36pt,dash pattern=on 1.35pt off 1.35pt] (1,-1) -- (4,-1) (1,-2) -- (4,-2) (1,-3) -- (4,-3) (1,-1) -- (1,-3) (2,-1) -- (2,-3) (3,-1) -- (3,-3) (4,-1) -- (4,-3);
    \filldraw[fill=kgray,draw=black,line width=0.54pt] (1,-1) circle (0.3);
    \draw[black,line width=0.54pt] (0.788,-0.788) -- (1.212,-1.212) (0.788,-1.212) -- (1.212,-0.788);
    \filldraw[fill=kgray,draw=black,line width=0.54pt] (2,-1) circle (0.3);
    \draw[black,line width=0.54pt] (1.788,-0.788) -- (2.212,-1.212) (1.788,-1.212) -- (2.212,-0.788);
    \filldraw[fill=kgray,draw=black,line width=0.54pt] (3,-1) circle (0.3);
    \draw[black,line width=0.54pt] (2.788,-0.788) -- (3.212,-1.212) (2.788,-1.212) -- (3.212,-0.788);
    \filldraw[fill=kgray,draw=black,line width=0.54pt] (4,-1) circle (0.3);
    \draw[black,line width=0.54pt] (3.788,-0.788) -- (4.212,-1.212) (3.788,-1.212) -- (4.212,-0.788);
    \filldraw[fill=kgray,draw=black,line width=0.54pt] (1,-2) circle (0.3);
    \draw[black,line width=0.54pt] (0.788,-1.788) -- (1.212,-2.212) (0.788,-2.212) -- (1.212,-1.788);
    \filldraw[fill=white,draw=black,line width=0.54pt] (2,-2) circle (0.3);
    \filldraw[fill=white,draw=black,line width=0.54pt] (3,-2) circle (0.3);
    \filldraw[fill=white,draw=black,line width=0.54pt] (4,-2) circle (0.3);
    \filldraw[fill=kgray,draw=black,line width=0.54pt] (1,-3) circle (0.3);
    \draw[black,line width=0.54pt] (0.788,-2.788) -- (1.212,-3.212) (0.788,-3.212) -- (1.212,-2.788);
    \filldraw[fill=white,draw=black,line width=0.54pt] (2,-3) circle (0.3);
    \filldraw[fill=white,draw=black,line width=0.54pt] (3,-3) circle (0.3);
    \filldraw[fill=white,draw=black,line width=0.54pt] (4,-3) circle (0.3);
    \node[lab] at (1,-0.45) {1};
    \node[lab] at (2,-0.45) {2};
    \node[lab] at (3,-0.45) {3};
    \node[lab] at (4,-0.45) {4};
    \node[lab] at (0.45,-1) {1};
    \node[lab] at (0.45,-2) {2};
    \node[lab] at (0.45,-3) {3};
    \node[capt,anchor=north] at (2.50,-3.45) {(b) our blackout};
  \end{scope}
  \begin{scope}[xshift=9.360cm]
    \draw[black,line width=0.36pt,dash pattern=on 1.35pt off 1.35pt] (1,-1) -- (4,-1) (1,-2) -- (4,-2) (1,-3) -- (4,-3) (1,-1) -- (1,-3) (2,-1) -- (2,-3) (3,-1) -- (3,-3) (4,-1) -- (4,-3);
    \draw[kred,line width=2.55pt,line join=round] (2,-2) -- (3,-1) -- (4,-2) -- (3,-3) -- cycle;
    \filldraw[fill=kgray,draw=black,line width=0.54pt] (1,-1) circle (0.3);
    \draw[black,line width=0.54pt] (0.788,-0.788) -- (1.212,-1.212) (0.788,-1.212) -- (1.212,-0.788);
    \filldraw[fill=kgray,draw=black,line width=0.54pt] (2,-1) circle (0.3);
    \draw[black,line width=0.54pt] (1.788,-0.788) -- (2.212,-1.212) (1.788,-1.212) -- (2.212,-0.788);
    \filldraw[fill=kgray,draw=black,line width=0.54pt] (3,-1) circle (0.3);
    \draw[black,line width=0.54pt] (2.788,-0.788) -- (3.212,-1.212) (2.788,-1.212) -- (3.212,-0.788);
    \filldraw[fill=kgray,draw=black,line width=0.54pt] (4,-1) circle (0.3);
    \draw[black,line width=0.54pt] (3.788,-0.788) -- (4.212,-1.212) (3.788,-1.212) -- (4.212,-0.788);
    \filldraw[fill=kgray,draw=black,line width=0.54pt] (1,-2) circle (0.3);
    \draw[black,line width=0.54pt] (0.788,-1.788) -- (1.212,-2.212) (0.788,-2.212) -- (1.212,-1.788);
    \filldraw[fill=white,draw=black,line width=0.54pt] (2,-2) circle (0.3);
    \filldraw[fill=white,draw=black,line width=0.54pt] (3,-2) circle (0.3);
    \filldraw[fill=white,draw=black,line width=0.54pt] (4,-2) circle (0.3);
    \filldraw[fill=kgray,draw=black,line width=0.54pt] (1,-3) circle (0.3);
    \draw[black,line width=0.54pt] (0.788,-2.788) -- (1.212,-3.212) (0.788,-3.212) -- (1.212,-2.788);
    \filldraw[fill=white,draw=black,line width=0.54pt] (2,-3) circle (0.3);
    \filldraw[fill=white,draw=black,line width=0.54pt] (3,-3) circle (0.3);
    \filldraw[fill=white,draw=black,line width=0.54pt] (4,-3) circle (0.3);
    \node[lab] at (1,-0.45) {1};
    \node[lab] at (2,-0.45) {2};
    \node[lab] at (3,-0.45) {3};
    \node[lab] at (4,-0.45) {4};
    \node[lab] at (0.45,-1) {1};
    \node[lab] at (0.45,-2) {2};
    \node[lab] at (0.45,-3) {3};
    \node[capt,anchor=north] at (2.50,-3.45) {(c) the rectangle};
  \end{scope}
  \begin{scope}[xshift=14.040cm]
    \draw[black,line width=0.36pt,dash pattern=on 1.35pt off 1.35pt] (1,-1) -- (4,-1) (1,-2) -- (4,-2) (1,-3) -- (4,-3) (1,-1) -- (1,-3) (2,-1) -- (2,-3) (3,-1) -- (3,-3) (4,-1) -- (4,-3);
    \filldraw[fill=kyellow,draw=black,line width=0.54pt] (2,-2) circle (0.3);
    \filldraw[fill=kyellow,draw=black,line width=0.54pt] (4,-2) circle (0.3);
    \filldraw[fill=kyellow,draw=black,line width=0.54pt] (3,-3) circle (0.3);
    \node[lab] at (1,-0.45) {1};
    \node[lab] at (2,-0.45) {2};
    \node[lab] at (3,-0.45) {3};
    \node[lab] at (4,-0.45) {4};
    \node[lab] at (0.45,-1) {1};
    \node[lab] at (0.45,-2) {2};
    \node[lab] at (0.45,-3) {3};
    \node[capt,anchor=north] at (2.50,-3.45) {(d) what we are shown};
  \end{scope}
\end{tikzpicture}
}
\caption{One round on a largest valid blackout of the $3\times4$ grid. Points are
(column, row). Gray crossed circles are black. Panel (c) shows the red rectangle,
and (d) its yellow corners. The three yellow corners fix the fourth by
Lemma~\ref{lem:threecorners}.}
\label{fig:game}
\end{figure}

Write $f(h,m)$ for the largest size of a valid blackout, and $a_h(m)$ for the
number of valid blackouts of that size. We call them \emph{optima}. The
values $f(h,m)$ form the table \seqnum{A397315} of the OEIS~\cite{oeis}, and
the counts $a_h(m)$ form the table \seqnum{A399580}. The companion
paper~\cite{pemmasaniExtremal} is about $f$. It proves the bound
$f(h,m)\le m+\lfloor h^2/4\rfloor$ with its cases of equality, shows that
equality holds for all large $m$, finds the order of the least such width, and
gives $f(h,m)$ for every $h\le6$. This paper is about the counts. It uses from
the companion paper only the statements listed in Theorem~\ref{thm:imports},
and the companion paper uses nothing from this one.

Put $q_h=\lfloor h^2/4\rfloor$. We record a blackout by the set of black rows
of each column, read from left to right. The results are as follows.

\begin{enumerate}[(1),leftmargin=*]
\item \emph{Two rows.} We have $a_2(m)=m2^{m-1}$ for $m\ge3$ (\seqnum{A001787}).
Up to symmetry the optima are counted by \seqnum{A322284}, by a bijection with
placements of kings (Theorem~\ref{thm:strip} and Proposition~\ref{prop:kings}).
\item \emph{Three rows.} We find $a_3(m)$ in closed form
(Theorem~\ref{thm:threecount}). For $m\ge4$ it is $4(m-1)(3m+2)\,3^{m-5}+E_m$,
where $E_m$ counts the optima with a column black in rows $1$ and $3$ and has
order $m^2(1+\sqrt3)^m$. So $a_3(m)\sim\frac4{81}m^23^m$. These sequences are
\seqnum{A401223} and \seqnum{A398448}. The series $\sum_{m\ge1}a_3(m)z^m$ is
rational with denominator $(1-3z)^3(1-2z-2z^2)^3(1+2z^2)^3$
(Theorem~\ref{thm:threegf}). The optima up to symmetry are counted in
Theorem~\ref{thm:threeorbits}, which gives \seqnum{A401226}.
\item \emph{Four rows.} For $m\ge4$ no valid blackout hides all $4$ corners of a
rectangle (Theorem~\ref{thm:hidden}). For $m\ge6$ a finite automaton with $6$
columns of memory accepts exactly the optima
(Proposition~\ref{prop:fourautomaton}). With computer help we find the generating function of $a_4(m)$
(\seqnum{A401224}). It is rational with a reduced denominator of degree $420$, and
$a_4(m)\sim c\,m^4\rho_4^m$, where $\rho_4=3.35199\ldots$ is algebraic of
degree $20$ (Theorem~\ref{thm:fourgfexact} and Corollary~\ref{cor:fourgrowth}).
\item \emph{Five rows.} Two transfer counters give $a_5(m)$ for $8\le m\le23$
(\seqnum{A401225}, Table~\ref{tab:fivevalues}).
\item \emph{Fixed height.} For every $h$ the valid column words form a regular
language, $f(h,m)-m$ is ultimately periodic in $m$, and $\sum_{m\ge1}a_h(m)z^m$
is rational (Theorems~\ref{thm:regular} and~\ref{thm:rational}). The proof does
not use the companion paper.
\item \emph{Tilted rectangles.} A tilted lattice rectangle whose bounding box has
height $n$ has width at most $\lfloor n^2/4\rfloor+1$. This proves a conjecture
of Kagey in \seqnum{A338885}, that the largest entry of its row $n$ is
\seqnum{A033638}$(n)$ (Remark~\ref{cor:oeisconj}).
\end{enumerate}

Statements that rest on a computation are marked as computer-assisted, and
Section~\ref{sec:verify} names the programs and the stored evidence.

\paragraph{Related work.} The value $f(3,m)=m+2$ was also proved
in Lean by the Omega Institute~\cite{trureturing2026}. The row bound
$f(h,m)\le m+q_h$ of the companion paper is, in graph form, due to
Neuwirth~\cite[Theorem~1(ii) and Proposition~17]{neuwirth2001}. Counting
configurations on a grid of fixed height by a transfer matrix is a standard
method~\cite[Section~4.7]{stanley2012}, \cite[Proposition~V.6]{flajolet2009}.
Examples are the independent sets of grid
graphs~\cite{calkinwilf1998,euler2005}, the nonattacking kings of
Wilf~\cite{wilf1995}, and the dominating sets of grids and king graphs counted
by Mertens~\cite{mertens2024}.
